(a) Since $k$ is infinite, every finite-dimensional affine or projective space has cardinality $|k|$, giving $|X|\le|k|$. Choose a nonempty affine open $U\subseteq X$ of dimension $r>0$. Noether normalization gives a finite integral inclusion $k[t_1,\ldots,t_r]\subseteq A(U)$. Lying over and the Nullstellensatz imply that the resulting map $U\to\mathbb A^r$ is surjective on closed points. Hence $|X|\ge|\mathbb A^r|=|k|$.

(b) Every proper closed subset of a curve is finite, since its irreducible components have dimension zero. Thus curves have the cofinite topology. By (a), any two curves have the same cardinality, so any bijection between them is a homeomorphism.
